\subsection{诺特完备局部环的结构}

设 A 和 C 是完备诺特局部环，且 u 是从 C 到 A 的局部同态，它经过取商诱导出从 \(\kappa_{C}\) 到 \(\kappa_{A}\) 的同构。设 \((p_{1}, \ldots, p_{m})\) 是生成 C 的理想 \(m_{C}\) 的一个序列，并令 \(t_{1}, \ldots, t_{n}\) 为 \(m_{A}\) 中的元素。置 \(B = C[[T_{1}, \ldots, T_{n}]]\)。

引理 3．——a) 存在唯一同态 \(v: \mathbf{B} \to \mathbf{A}\)，它延拓 \(u\)，并将 \(\mathbf{T}_i\) 映到 \(t_i\)，对 \(1 \leqslant i \leqslant n\) 均如此。

\begin{enumerate}
\def\labelenumi{\alph{enumi})}
\setcounter{enumi}{1}
\item
  v 为满射的充要条件是序列 \((u(p_{1}), \ldots, u(p_{m}), t_{1}, \ldots, t_{n})\) 生成 A 的理想 \(m_{A}\)，或者等价地，这些元素模 \(m_{A}^{2}\) 的类生成 \(m_{A}/m_{A}^{2}\)，后者是域 \(k_{A}\) 上的向量空间。
\item
  要使 \(v\) 使 A 成为有限 B-代数，充要条件是序列 \((u(p_1), \ldots, u(p_m), t_1, \ldots, t_n)\) 生成 A 的 \(\mathfrak{m}_{\mathbf{A}}\)-进拓扑的定义理想。
\end{enumerate}

令 n 表示由 \(T_{1}, \ldots, T_{n}\) 生成的环 B 的理想。每个从 B 到 A、延拓 u 且满足 \(v(T_{i}) = t_{i}\) 的同态都将 n 映入 \(m_{A}\)，所以当 B 赋予 n-进拓扑时它是连续的。于是 v 的存在性与唯一性由 A，IV，p. 26，命题 4 得出。

环 \(B = C[[T_{1}, \ldots, T_{n}]]\) 是完备诺特局部环（III，§ 2，第 10 号，定理 2 的推论 6及第 6 号，命题 6），其极大理想 \(m_{B}\) 由 \(p_{1}, \ldots, p_{m}, T_{1}, \ldots, T_{n}\) 生成。因此有 \(v(m_{B}) \subset m_{A}\)，且 v 定义了从 gr(B) = \(\sum_{n=0}^{\infty} m_{B}^{n}/m_{B}^{n+1}\) 到 gr(A) = \(\sum_{n=0}^{\infty} m_{A}^{n}/m_{A}^{n+1}\) 的同态 gr(v)。现在环 gr(A) 由 \(A/m_{A} = \kappa_{A}\) 和 \(m_{A}/m_{A}^{2}\) 生成；gr(v) 诱导出从 \(\kappa_{B} = \kappa_{C}\) 到 \(\kappa_{A}\) 的同构，并且在模 \(m_{B}^{2}\) 后，元素 \(p_{1}, \ldots, p_{m}, T_{1}, \ldots, T_{n}\) 的类生成 \(m_{B}/m_{B}^{2}\)，作为 \(\kappa_{B}\) 上的向量空间；此外，v 是满射当且仅当 gr(v) 是满射（III，§ 2，第 8 号，定理 1 的推论 2）。这证明了 b)。

由序列 \((u(p_{1}), \ldots, u(p_{m}), t_{1}, \ldots, t_{n})\) 生成的 A 的理想正是 \(v(\mathfrak{m}_{\mathrm{B}})\)A。由于 \(m_{A}\) 包含 \(v(\mathfrak{m}_{\mathrm{B}})\)，A 关于 \(v(\mathfrak{m}_{\mathrm{B}})\)A-进拓扑是一个 Zariski 环。环 \(A/v(\mathfrak{m}_{\mathrm{B}})\)A 是阿廷环的充要条件是其作为 A-模的长度有限。但是由于 A 上每个简单模都被 \(m_{A}\) 湮灭，且根据假设 \(A/m_{A}\) 与 \(B/m_{B}\) 同构，这等价于域 \(B/m_{B}\) 上的向量空间 \(A/v(\mathfrak{m}_{\mathrm{B}})\)A 的维数有限。根据 IV，§ 2，第 5 号，命题 9 的推论 2，可知 \(v(\mathfrak{m}_{\mathrm{B}})\)A 是 A 的定义理想，当且仅当 \(A/v(\mathfrak{m}_{\mathrm{B}})\)A 在 \(B/m_{B}\) 上的维数有限。若 A 是有限 B-代数，确实如此。

假设 \(v(\mathfrak{m}_{\mathrm{B}})\)A 是 A 的定义理想。B-模 A 的 \(\mathfrak{m}_{\mathrm{B}}\)-进拓扑于是与环 A 的 \(\mathfrak{m}_{\mathrm{A}}\)-进拓扑重合，因而是分离的。由于 \(\mathrm{A} / v(\mathfrak{m}_{\mathrm{B}})\)A 是 \(\mathrm{B} / \mathfrak{m}_{\mathrm{B}}\) 上的有限生成模，A 是有限生成 B-模（III，§ 2，第 3 号，例 3 及第 9 号，命题 12 的推论 1）。这证明了 c)。

引理 4．——假设诺特局部环 C 是正则的，且 \((p_1, \ldots, p_m)\) 是 C 的一个坐标系（VIII，§ 5，第 1 号，定义 1）。

\begin{enumerate}
\def\labelenumi{\alph{enumi})}
\item
  若序列 \((u(p_1), \ldots, u(p_m), t_1, \ldots, t_n)\) 是 A 的一个正则序列（VIII，§ 3，第 2 号，定义 1），则同态 \(v: \mathbf{B} \to \mathbf{A}\) 是单射。
\item
  要使 \(v\) 为单射且使 A 成为 B 上的有限代数，充要条件是 \((u(p_1), \ldots, u(p_m), t_1, \ldots, t_n)\) 是 A 的极大正则序列。此时 A 的维数为 \(m + n\)。
\end{enumerate}

序列 \((u(p_1), \ldots, u(p_m), t_1, \ldots, t_n)\) 成为 A 的极大正则序列的充要条件是：它生成 A 的定义理想，且 A 的维数为 \(m + n\)（VIII，§ 3，第 2 号，定理 1）。根据引理 3 c)，这等价于 A 是有限 B-代数且维数为 \(m + n\)。现在 C 是维数为 \(m\) 的整诺特环，从而 \(\mathbf{B} = \mathbf{C}[[\mathbf{T}_1, \ldots, \mathbf{T}_n]]\) 是维数为 \(m + n\) 的整诺特环（VIII，§ 3，第 4 号，命题 8 的推论 3）。若 A 是有限 B-代数，且 \(\mathfrak{a}\) 是 \(v\) 的核，则有 \(\dim(\mathbf{A}) = \dim(\mathbf{B}/\mathfrak{a})\)（VIII，§ 2，第 3 号，定理 1，c)）；由于 B 是有限维整环，有 \(\dim(\mathbf{B}/\mathfrak{a}) < \dim(\mathbf{B})\)，当 \(\mathfrak{a} \neq \{0\}\) 时（VIII，§ 1，第 3 号，命题 6，e)）。因此，若 A 是有限 B-代数，则 \(v\) 是单射当且仅当 A 的维数为 \(m + n\)。这证明了 b)。

假设序列 \((u(p_1), \ldots, u(p_m), t_1, \ldots, t_n)\) 是 \(\mathfrak{m}_{\mathbf{A}}\) 中的正则元素序列。根据（VIII，§ 3，第 2 号，定理 1），可以在其中添入元素 \(t_{n+1}, \ldots, t_{n+r}\)，这些元素属于 \(\mathfrak{m}_{\mathbf{A}}\)，使其成为极大正则序列。根据上面的结果，存在单射同态 \(w\)，其从 \(\mathrm{C}[[\mathrm{T}_1, \ldots, \mathrm{T}_n, \mathrm{T}_{n+1}, \ldots, \mathrm{T}_{n+r}]] = \mathrm{B}[[\mathrm{T}_{n+1}, \ldots, \mathrm{T}_{n+r}]]\) 延拓 \(v\)，并将 \(\mathrm{T}_{n+j}\) 映到 \(t_{n+j}\)，对 \(1 \leqslant j \leqslant r\) 均如此。因此 \(v\) 是单射。这证明了 \(a\)。

定理 3．——设 A 是局部诺特完备环，其剩余域 k 的特征为 p。设 C 是无限长度的 p-环，其剩余域同构于 k（第 3 号，命题 5）。

\begin{enumerate}
\def\labelenumi{\alph{enumi})}
\item
  令 m 为向量空间 \(\mathfrak{m}_{\mathbf{A}}/(\mathfrak{m}_{\mathbf{A}}^{2} + p\mathbf{A})\) 在域 k 上的维数。存在环 \(\mathrm{C}[[\mathrm{T}_{1}, ..., \mathrm{T}_{m}]]\) 的理想 a，使 A 同构于 \(\mathrm{C}[[\mathrm{T}_{1}, ..., \mathrm{T}_{m}]]/\mathfrak{a}\)。
\item
  令 d 为 A 的维数。假设 \(p1_{A}\) 不是 A 中的零因子。则存在 A 的子环 A'，同构于 \(C[[T_{1}, \ldots, T_{d-1}]]\)，且 A 是 A' 上的有限代数。
\end{enumerate}

令 C' 为 A 的一个科恩子环（第 2 号，定理 1）。由于 C 的长度无限，根据第 3 号的命题 4，存在从 C 到 C' 的同态。因此存在局部同态 \(u:C\to A\)。选取元素 \(t_{1},...,t_{m}\)，这些元素属于 \(\mathfrak{m}_{\mathbf{A}}\)，其类构成向量空间 \(\mathfrak{m}_{\mathbf{A}}/(\mathfrak{m}_{\mathbf{A}}^{2}+p\mathbf{A})\) 在域 \(k\) 上的一个基。有 \(u(p1_{\mathbf{C}})=p1_{\mathbf{A}}\)，且引理 3 的 \(b)\) 证明了存在从 \(\mathrm{C}[[\mathrm{T}_{1},..., \mathrm{T}_{m}]]\) 到 A 的满射同态，它延拓 \(u\)，并将 \(\mathrm{T}_{i}\) 映到 \(t_{i}\)，对 \(1\leqslant i\leqslant m\) 均如此。这证明了 \(a)\)。

若 \(p1_{\mathbf{A}}\) 不是 A 中的零因子，因而是 A 的正则元素（VIII，§ 3，第 2 号，命题 3）。于是存在（VIII，§ 3，第 2 号，定理 1）元素 \(t_1, \ldots, t_{d-1}\)，这些元素属于 \(\mathfrak{m}_{\mathbf{A}}\)，使序列 \((p1_{\mathbf{A}}, t_1, \ldots, t_{d-1})\) 成为 A 的极大正则序列。诺特局部环 C 是正则的，且 \((p1_{\mathbf{C}})\) 是 C 的一个坐标系。定理 3 的断言 \(b)\) 随后由引理 4 的 \(b)\) 得出。

